\documentclass[10pt]{article}
\usepackage[T1]{fontenc}
\usepackage[utf8]{inputenc}
\usepackage[french]{babel}
\usepackage{amsmath,amssymb,mathrsfs,mathpazo}
\usepackage{geometry}
\geometry{a4paper,margin=2cm}
\usepackage{graphicx,array,tabularx,tabularray,colortbl,multirow,multicol,booktabs,enumitem,pifont,ulem,fancybox,diagbox,cancel}
\usepackage{pstricks,pstricks-add,pst-plot,pst-tree,pst-node,pst-text,pst-func,pst-math,pst-bspline,pst-3dplot,pst-eucl,pst-solides3d}
\let\pscalculate\undefined
\usepackage{pst-bezier}
\usepackage{tikz,pgfplots,tkz-tab}
\usetikzlibrary{arrows,arrows.meta,patterns,calc,positioning,trees,decorations.markings,shapes.geometric,shapes.misc,shapes.arrows,intersections,decorations.pathmorphing,decorations.pathreplacing,backgrounds,shadows,babel}
\pgfplotsset{compat=1.18}
\usepgfplotslibrary{fillbetween}
\usepackage[np]{numprint}
\usepackage{listings,makecell,pas-tableur,twoopt}
\let\barre\undefined
\usepackage[pstricks]{tablvar}
\definecolor{darkgreen}{RGB}{0,100,0}
\providecommand{\red}{\color{red}}
\providecommand{\blue}{\color{blue}}
\providecommand{\green}{\color{green}}
\providecommand{\black}{\color{black}}
\providecommand{\white}{\color{white}}
\providecommand{\gray}{\color{gray}}
\providecommand{\cyan}{\color{cyan}}
\providecommand{\magenta}{\color{magenta}}
\providecommand{\yellow}{\color{yellow}}
\providecommand{\violet}{\color{violet}}
\providecommand{\orange}{\color{orange}}
\providecommand{\euro}{\texteuro}
\providecommand{\EUR}{\texteuro}
\providecommand{\ieme}{\textsuperscript{e}}
\providecommand{\degres}{\ensuremath{^\circ}}
\providecommand{\degree}{\ensuremath{^\circ}}
\providecommand{\vv}[1]{\overrightarrow{#1}}
\providecommand{\Coord}[1]{#1}
\providecommand{\decosix}{}
\providecommand{\starredbullet}{\ensuremath{\bullet}}
\pagestyle{empty}
\setlength{\parindent}{0pt}
\usepackage{adjustbox}\DeclareUnicodeCharacter{2236}{:}\DeclareUnicodeCharacter{2190}{\ensuremath{\leftarrow}}
\providecommand{\R}{\mathbb{R}}
\providecommand{\N}{\mathbb{N}}
\providecommand{\Z}{\mathbb{Z}}
\providecommand{\Q}{\mathbb{Q}}
\providecommand{\D}{\mathbb{D}}
\providecommand{\C}{\mathbb{C}}
\providecommand{\e}{{\mathrm{e}}}
\providecommand{\vect}[1]{\overrightarrow{#1}}
\providecommand{\vectt}[1]{\overrightarrow{\mathrm{#1}}}
\providecommand{\vvt}[1]{\overrightarrow{\mathrm{#1}}}
\providecommand{\np}[2][]{{#2~#1}}
\providecommand{\vv}[1]{{\overrightarrow{#1}}}
\renewcommand{\d}{{\mathrm{d}}}
\providecommand{\degres}{{^\circ}}
\providecommand{\degre}{{^\circ}}
\providecommand{\blue}{\color{blue}}
\providecommand{\red}{\color{red}}
\providecommand{\enskip}{\quad}
\providecommand{\llbracket}{\lbrack\!\lbrack}
\providecommand{\rrbracket}{\rbrack\!\rbrack}
\providecommand{\euro}{{\text{€}}}
\providecommand{\no}{{\text{n°}}}
\providecommand{\interval}[2]{{[#1;#2]}}
\providecommand{\intervFO}[2]{{[#1;#2[}}
\providecommand{\dotfill}{\ldots}

\providecommand{\ds}{\displaystyle}
\providecommand{\textding}[1]{#1}
\providecommand{\Oij}{\(\mathrm{O};\vec{\imath},\vec{\jmath}\)}
\providecommand{\Oijk}{\(\mathrm{O};\vec{\imath},\vec{\jmath},\vec{k}\)}
\providecommand{\Ouv}{\(\mathrm{O};\vec{u},\vec{v}\)}
\providecommand{\intervFF}[2]{\ensuremath{[#1;#2]}}
\renewcommand{\euro}{\ifmmode\text{\texteuro}\else\texteuro\fi}
\DeclareUnicodeCharacter{00B0}{\ifmmode^{\circ}\else\textdegree\fi}
\DeclareUnicodeCharacter{20AC}{\ifmmode\text{\texteuro}\else\texteuro\fi}
% Énoncés : Éducation nationale. Transcription et corrigés : APMEP, Denis Vergès et contributeurs. Découpage et classement : CoopMaths, Sébastien Lozano, Éric Elter et contributeurs. Adaptation de la présentation : Maths Explorer.
% CC BY-SA : https://coopmaths.fr/www/about/
% https://coopmaths.fr/alea/static/bac/2025/tex/bac_2025_06_sujet1_asie_1_cor.tex
\begin{document}
\noindent\textit{APMEP / CoopMaths. CC BY-SA. Présentation adaptée par Maths Explorer.}\par\medskip
\medskip

\textbf{Affirmation 1 : Fausse.}

Le vecteur $ \vect{\mathrm{AB}} $ a pour coordonnées $ \vect{\mathrm{AB}}\begin{pmatrix} 1 \\ 0 \\ 0\end{pmatrix} $ et $ \vect{\mathrm{AC}} $ a pour coordonnées : $ \vect{\mathrm{AC}}\begin{pmatrix} \alpha - 1 \\2\\\alpha \end{pmatrix} $.

Quel que soit le réel $\alpha$, ces coordonnées ne sont pas proportionnelles, donc les vecteurs ne sont pas colinéaires, et ils définissent bien un plan.

Cependant, le repère étant orthonormé, on peut utiliser les coordonnées pour calculer un produit scalaire : \quad
$ \vect{\mathrm{AC}}\cdot\vect{\jmath} = (\alpha-1)\times 0 + 2 \times 1 + \alpha \times 0 = 2$.

On constate que les vecteurs $ \vect{\jmath} $ et $ \vect{\mathrm{AC}} $ ne sont pas orthogonaux, donc $ \vect{\jmath} $ n'est pas un vecteur normal au plan (ABC) : cette partie de l'affirmation est donc fausse.

\medskip

\textbf{Affirmation 2 : Fausse.}

En interprétant le système de représentation paramétrique, la droite $(d)$ est dirigée par $\vect{u} \begin{pmatrix}1 \\ 2  \\ -1 \end{pmatrix}$. On rappelle : $ \vect{\mathrm{AC}}\begin{pmatrix} \alpha - 1 \\2\\\alpha \end{pmatrix} $.

Pour que les droites soient parallèles, il faudrait que les vecteurs $ \vect{u} $ et $ \vect{\mathrm{AC}} $ soient proportionnels, or, ici, ils ont la même ordonnée, donc ils devraient être égaux.

Or, si on choisit $\alpha$ égal à 1, les altitudes seront égales, mais pas les abscisses, si on choisit $\alpha$ égal à 2, alors les abscisses seront égales, mais pas les altitudes, et, pour tout autre réel $\alpha$, les abscisses et altitudes seront différentes entre les deux vecteurs.

Quel que soit le réel $\alpha$, les vecteurs $ \vect{u} $ et $ \vect{\mathrm{AC}} $ sont donc non colinéaires, et les droites ne sont pas parallèles : l'affirmation est donc fausse.

\medskip

\textbf{Affirmation 3 : Vraie.}

O étant l'origine du repère, les coordonnées des vecteurs $ \vect{\mathrm{OA}} $ sont celles du point A, donc celles de $ \vect{\mathrm{AO}} $ sont : $ \begin{pmatrix}-1\\-1\\0\end{pmatrix}$. On rappelle $ \vect{\mathrm{AB}}\begin{pmatrix} 1 \\ 0 \\ 0\end{pmatrix} $.

On a donc, dans le repère orthonormé : \quad$ \vect{\mathrm{AO}} \cdot \vect{\mathrm{AB}} = -1 \times 1 + (-1) \times 0 + 0 \times 0 = -1$.

Par ailleurs, on a : \quad $ \mathrm{AO} = \sqrt{(-1)^2 + (-1)^2 + 0^2} = \sqrt{2} $, \quad et \quad $ \mathrm{AB} = \sqrt{1^2 + 0^2 + 0^2} = \sqrt{1} = 1$.

On a donc aussi : \quad $ \vect{\mathrm{AO}} \cdot \vect{\mathrm{AB}} = \mathrm{AO} \times \mathrm{AB} \times \cos \left( \widehat{\mathrm{OAB}}\right) = \sqrt{2} \cos \left( \widehat{\mathrm{OAB}}\right)$.

En égalant les deux expressions du produit scalaire, il vient :
\quad $-1 = \sqrt{2}\cos \left(\widehat{\mathrm{OAB}}\right)$,

ce qui équivaut à :\quad $\cos \left(\widehat{\mathrm{OAB}}\right) = \dfrac{-1}{\sqrt{2}} = \dfrac{-\sqrt{2}}{2}$.

On reconnaît ici l'opposé du cosinus d'un angle de $ \dfrac{\pi}{4} $, ce qui est donc le cosinus d'un angle de mesure $\pi - \dfrac{\pi}{4} = \dfrac{3\pi}{4}$, soit effectivement 135\ensuremath{^\circ}.

\medskip

\textbf{Affirmation 4 : Fausse.}

Avec les coordonnées données pour H et la représentation paramétrique de $(d)$, pour avoir $x_\mathrm{H} = 1 + t$, il faut avoir $t = 0$, et donc les ordonnées et altitudes ne correspondront pas : le point H n'est pas sur la droite $(d)$, donc il n'est pas le projeté orthogonal de A sur $(d)$.

\medskip


\textbf{Affirmation 5 : Vraie}

Soit $t$ un réel quelconque et $M_t$ le point de paramètre $t$ sur la droite $(d)$.

La distance $\mathrm{O}M_t$ est donnée par :

$\mathrm{O}M_t = \sqrt{(1 + t)^2 + (2t)^2 + (-t)^2} = \sqrt{1 + 2t + t^2 + 4t^2 + t^2} = \sqrt{6t^2 + 2t + 1} $.

On a donc : \quad $ \begin{aligned}[t] \mathrm{O}M_t = 1
	&\iff \sqrt{6t^2 + 2t + 1} = 1\\
	&\iff 6t^2 + 2t + 1 = 1 \\%\quad \text{car une distance est toujours positive}\\
	&\iff 6t^2 + 2t = 0\\
	&\iff 2t(3t + 1) = 0\\
	&\iff 2t = 0 \quad \text{ou} \quad 3t + 1 = 0\\
	&\iff t = 0 \quad \text{ou} \quad t = -\dfrac{1}{3}\\
\end{aligned} $

L'équation a deux solutions, il y a donc exactement deux points de la droite $(d)$ qui sont sur la sphère : les points de paramètre 0 et $ -\dfrac{1}{3} $, l'affirmation est vraie.
\medskip

\end{document}
